\documentclass[12pt,a4paper]{article}

% --- Metadata (per course) ---
\def\courseshortheader{XÁC SUẤT VÀ ĐỘ ĐO}
\def\coverbookline{NHÓM 5 --- XÁC SUẤT VÀ THỐNG KÊ}
\def\covermaintitle{Xác suất và độ đo}
\def\coversubtitle{Giáo án lý thuyết --- Không gian xác suất, hội tụ, kỳ vọng có điều kiện}

% Shared preamble for nhom_5 (requires \courseshortheader before \input)
\usepackage[utf8]{inputenc}
\usepackage[vietnamese]{babel}
\usepackage{amsmath,amssymb,amsthm}
\usepackage{geometry}
\usepackage{enumitem}
\usepackage[most]{tcolorbox}
\usepackage{hyperref}
\usepackage{fancyhdr}
\usepackage{graphicx}
\usepackage{booktabs}
\usepackage{tikz}
\usetikzlibrary{calc,arrows.meta,decorations.pathreplacing,positioning}
\usepackage{eso-pic}
\usepackage{xcolor}

\geometry{a4paper, margin=2cm, bottom=2.5cm}
\hypersetup{colorlinks=true, linkcolor=blue!70!black, urlcolor=blue}
\setlist[itemize]{topsep=4pt, itemsep=2pt}
\setlist[enumerate]{topsep=4pt, itemsep=2pt}

\AddToShipoutPictureFG{%
  \AtPageCenter{%
    \begin{tikzpicture}[remember picture, overlay]
      \node[rotate=45, scale=1, opacity=0.06]{%
        \fontsize{60}{72}\selectfont\bfseries\sffamily Đặng Tấn Quí%
      };
    \end{tikzpicture}%
  }%
}

\pagestyle{fancy}
\fancyhf{}
\fancyhead[L]{\small\textit{\courseshortheader}}
\fancyhead[R]{\small\textit{Đặng Tấn Quí}}
\fancyfoot[C]{\thepage}
\fancyfoot[L]{\footnotesize\textcopyright\ Đặng Tấn Quí}
\fancyfoot[R]{\footnotesize SĐT: 0377\,122\,210}
\renewcommand{\headrulewidth}{0.4pt}
\renewcommand{\footrulewidth}{0.4pt}

\fancypagestyle{plain}{%
  \fancyhf{}
  \fancyfoot[C]{\thepage}
  \fancyfoot[L]{\footnotesize\textcopyright\ Đặng Tấn Quí}
  \fancyfoot[R]{\footnotesize SĐT: 0377\,122\,210}
  \renewcommand{\headrulewidth}{0pt}
  \renewcommand{\footrulewidth}{0.4pt}
}

\newtcolorbox{dongco}[1][]{%
  enhanced, breakable,
  colback=teal!4!white, colframe=teal!60!black,
  fonttitle=\bfseries, title={Động cơ học -- Tại sao cần khái niệm này?}, #1%
}
\newtcolorbox{ynghia}[1][]{%
  enhanced, breakable,
  colback=purple!3!white, colframe=purple!50!black,
  fonttitle=\bfseries, title={Ý nghĩa \& Trực giác}, #1%
}
\newtcolorbox{ungdung}[1][]{%
  enhanced, breakable,
  colback=olive!5!white, colframe=olive!60!black,
  fonttitle=\bfseries, title={Ứng dụng thực tế}, #1%
}
\newtcolorbox{dinhnghia}[1][]{%
  enhanced, breakable,
  colback=blue!3!white, colframe=blue!60!black,
  fonttitle=\bfseries, title={Định nghĩa}, #1%
}
\newtcolorbox{dinhly}[1][]{%
  enhanced, breakable,
  colback=green!3!white, colframe=green!50!black,
  fonttitle=\bfseries, title={Định lý}, #1%
}
\newtcolorbox{tinhchat}[1][]{%
  enhanced, breakable,
  colback=yellow!5!white, colframe=orange!50!black,
  fonttitle=\bfseries, title={Tính chất}, #1%
}
\newtcolorbox{phuongphap}[1][]{%
  enhanced, breakable,
  colback=gray!5!white, colframe=gray!50!black,
  fonttitle=\bfseries, title={Phương pháp}, #1%
}
\newtcolorbox{luuy}[1][]{%
  enhanced, breakable,
  colback=red!3!white, colframe=red!50!black,
  fonttitle=\bfseries, title={Lưu ý}, #1%
}
\newtcolorbox{congthuc}[1][]{%
  enhanced, breakable,
  colback=cyan!3!white, colframe=cyan!50!black,
  fonttitle=\bfseries, title={Công thức}, #1%
}
\newtcolorbox{vidu}[1][]{%
  enhanced, breakable,
  colback=violet!3!white, colframe=violet!50!black,
  fonttitle=\bfseries, title={Ví dụ}, #1%
}
\newtcolorbox{phanvidu}[1][]{%
  enhanced, breakable,
  colback=red!2!white, colframe=red!40!black,
  fonttitle=\bfseries, title={Phản ví dụ}, #1%
}
\newtcolorbox{tongket}[1][]{%
  enhanced, breakable,
  colback=blue!4!white, colframe=blue!70!black,
  fonttitle=\bfseries, title={Tổng kết chương}, #1%
}


\begin{document}

\thispagestyle{empty}
\begin{tikzpicture}[remember picture, overlay]
  \fill[blue!80!black] (current page.south west) rectangle (current page.north east);
  \fill[blue!60!cyan, opacity=0.8]
    (current page.south west) -- (current page.north west) --
    ([xshift=-3cm]current page.north east) -- (current page.south east) -- cycle;
  \foreach \r/\o in {6/0.05, 4.5/0.08, 3/0.12} {
    \fill[white, opacity=\o] ([xshift=2cm, yshift=-2cm]current page.north east) circle (\r cm);
  }
  \draw[white, line width=2pt] ([yshift=-5cm]current page.north west) ++(2cm,0) -- ++(17cm,0);
  \draw[white, line width=2pt] ([yshift=-16cm]current page.north west) ++(2cm,0) -- ++(17cm,0);
  \node[anchor=west, white, font=\fontsize{28}{34}\bfseries\selectfont]
    at ([xshift=3cm, yshift=-7.5cm]current page.north west) {\coverbookline};
  \node[anchor=west, white, font=\fontsize{20}{26}\selectfont]
    at ([xshift=3cm, yshift=-9cm]current page.north west) {\covermaintitle};
  \node[anchor=west, white, font=\fontsize{16}{22}\selectfont]
    at ([xshift=3cm, yshift=-10.2cm]current page.north west) {\coversubtitle};
  \node[anchor=west, white, font=\Large\bfseries]
    at ([xshift=3cm, yshift=-13cm]current page.north west) {Biên soạn:};
  \node[anchor=west, yellow!80!white, font=\fontsize{24}{30}\bfseries\selectfont]
    at ([xshift=3cm, yshift=-14.5cm]current page.north west) {ĐẶNG TẤN QUÍ};
  \node[anchor=west, white!90, font=\large]
    at ([xshift=3cm, yshift=-18cm]current page.north west) {SĐT: 0377\,122\,210};
  \node[anchor=south, white!80, font=\small]
    at ([yshift=1.5cm]current page.south) {Tài liệu có bản quyền --- Cấm sao chép khi chưa được sự đồng ý của tác giả};
  \node[anchor=south east, white, font=\Large\bfseries]
    at ([xshift=-2cm, yshift=3cm]current page.south east) {\the\year};
\end{tikzpicture}

\newpage
\tableofcontents
\newpage

\section*{Lời nói đầu}
\addcontentsline{toc}{section}{Lời nói đầu}

Tài liệu thuộc \textbf{Nhóm 5 --- Xác suất và thống kê}, biên soạn theo cùng triết lý với giáo án Đại số tuyến tính: học để \emph{hiểu} và \emph{dùng}, không chỉ để ghi nhớ công thức.

\begin{enumerate}
  \item \textbf{Động cơ trước định nghĩa.} Mỗi phần lớn mở bằng ô \textsf{\color{teal!60!black} Động cơ học} --- câu hỏi thực tế hoặc bài toán mẫu cần khái niệm đó.
  \item \textbf{Trực giác trước công thức.} Ô \textsf{\color{purple!50!black} Ý nghĩa \& Trực giác} giải thích ``công thức đang nói gì'' (xác suất = tần suất dài hạn, kỳ vọng = trung bình có trọng số, v.v.).
  \item \textbf{Ví dụ có kiểm tra.} Ô \textsf{\color{violet!50!black} Ví dụ} nên tự làm trước khi đọc lời giải; phần thống kê ứng dụng cần gắn dữ liệu số cụ thể.
  \item \textbf{Tổng kết cuối mỗi chương.} Ô \textsf{\color{blue!70!black} Tổng kết chương} liệt kê kết quả phải nhớ và liên kết sang phần sau.
  \item \textbf{Phân tầng theo môn.} X1--X3 là nền; X4--X5 nâng cao lý thuyết; X6--X8 chuyên sâu (đa biến, Bayes, quá trình ngẫu nhiên).
\end{enumerate}

\noindent\textbf{Cách đọc:} Động cơ $\to$ Định nghĩa / Định lý $\to$ Ví dụ $\to$ Ý nghĩa $\to$ Tổng kết. Với môn thống kê, luôn hỏi: \emph{mẫu là gì, tham số nào chưa biết, giả thuyết $H_0$ là gì}.

\newpage


\section{Tại sao học xác suất đo lường?}

\begin{dongco}[title={Động cơ --- Tại sao học xác suất đo lường?}]
X1 cho trực giác; X4 xây nền \emph{chặt}: $\sigma$-đại số, các kiểu hội tụ ($P$, a.s., phân phối), hàm đặc trưng, kỳ vọng có điều kiện --- cần cho chứng minh luật số lớn và môn quá trình ngẫu nhiên.
\end{dongco}

\newpage

\section{Không gian xác suất}

\begin{dongco}[title={Động cơ --- Không gian xác suất}]
Ta cần khung toán cho phép xử lý vô hạn biến cố và tính xác suất nhất quán.
\end{dongco}

\begin{dinhnghia}
  $(\Omega, \mathcal{F}, P)$: $\Omega$ không gian mẫu, $\mathcal{F}$ là $\sigma$-đại số trên $\Omega$, $P: \mathcal{F} \to [0, 1]$ là độ đo xác suất ($P(\Omega) = 1$).
\end{dinhnghia}

\begin{dinhnghia}[title={Biến ngẫu nhiên (đo được)}]
  $X: (\Omega, \mathcal{F}) \to (\mathbb{R}, \mathcal{B})$ đo được: $X^{-1}(B) \in \mathcal{F}$ $\forall B \in \mathcal{B}(\mathbb{R})$.

  \textbf{Phân phối} (luật) của $X$: $P_X = P \circ X^{-1}$, là độ đo xác suất trên $(\mathbb{R}, \mathcal{B})$.
\end{dinhnghia}

\begin{dinhnghia}[title={Tích phân Lebesgue là kỳ vọng}]
  $E[X] = \displaystyle\int_\Omega X(\omega)\,dP(\omega) = \int_{\mathbb{R}} x\,dP_X(x)$.

  Tồn tại khi $E[|X|] < \infty$ ($X \in L^1(\Omega, P)$).
\end{dinhnghia}

%% ================================================================
%%  CHƯƠNG 2: CÁC LOẠI HỘI TỤ
%% ================================================================

\begin{tongket}[title={Tổng kết: Không gian xác suất}]
\begin{enumerate}
  \item Nắm lại các \textbf{định nghĩa} và \textbf{công thức} cốt lõi trong mục ``Không gian xác suất''.
  \item Tự làm lại ít nhất một ví dụ (nếu có) hoặc áp dụng công thức vào một tình huống số.
  \item Ghi chú điều kiện áp dụng (giả thiết độc lập, phân phối chuẩn, $n$ đủ lớn, v.v.).
\end{enumerate}
\end{tongket}

\section{Các loại hội tụ}

\begin{dongco}[title={Động cơ --- Các loại hội tụ}]
Dãy BNN và dãy phân phối ``tiến tới đâu'' --- cần phân biệt hội tụ theo phân phối, theo xác suất, hầu chắc.
\end{dongco}

\begin{dinhnghia}
  Cho $(X_n)$ dãy BNN:
  \begin{itemize}
    \item \textbf{Hội tụ hầu chắc chắn} (a.s.): $P\!\left(\lim_{n \to \infty} X_n = X\right) = 1$, viết $X_n \xrightarrow{\text{a.s.}} X$.
    \item \textbf{Hội tụ theo xác suất:} $\forall \varepsilon > 0$: $P(|X_n - X| > \varepsilon) \to 0$, viết $X_n \xrightarrow{P} X$.
    \item \textbf{Hội tụ trong $L^p$:} $E[|X_n - X|^p] \to 0$, viết $X_n \xrightarrow{L^p} X$.
    \item \textbf{Hội tụ theo phân phối:} $F_{X_n}(x) \to F_X(x)$ tại mọi điểm liên tục của $F_X$, viết $X_n \xrightarrow{d} X$.
  \end{itemize}
\end{dinhnghia}

\begin{dinhly}[title={Quan hệ giữa các loại hội tụ}]
  \[
    \text{a.s.} \Rightarrow P \Rightarrow d, \qquad L^p \Rightarrow P \Rightarrow d \quad(p \ge 1).
  \]
  Ngược lại không đúng nói chung, ngoại trừ:
  \begin{itemize}
    \item $X_n \xrightarrow{P} X$ $\Rightarrow$ $\exists$ dãy con $X_{n_k} \xrightarrow{\text{a.s.}} X$.
    \item $X_n \xrightarrow{d} c$ (hằng số) $\Leftrightarrow$ $X_n \xrightarrow{P} c$.
  \end{itemize}
\end{dinhly}

\begin{dinhly}[title={Bổ đề Borel--Cantelli}]
  \begin{enumerate}
    \item Nếu $\sum P(A_n) < \infty$ thì $P(\limsup A_n) = 0$ (``hầu hết hữu hạn $A_n$ xảy ra'').
    \item Nếu $A_n$ \textbf{độc lập} và $\sum P(A_n) = \infty$ thì $P(\limsup A_n) = 1$.
  \end{enumerate}
\end{dinhly}

\begin{dinhly}[title={Định lý Slutsky}]
  Nếu $X_n \xrightarrow{d} X$ và $Y_n \xrightarrow{P} c$ (hằng) thì:
  $X_n + Y_n \xrightarrow{d} X + c$, \;$X_n Y_n \xrightarrow{d} cX$, \;$X_n/Y_n \xrightarrow{d} X/c$ ($c \ne 0$).
\end{dinhly}

\begin{dinhly}[title={Định lý ánh xạ liên tục (Continuous Mapping)}]
  Nếu $X_n \xrightarrow{d} X$ và $g$ liên tục thì $g(X_n) \xrightarrow{d} g(X)$. Tương tự cho a.s. và theo $P$.
\end{dinhly}

%% ================================================================
%%  CHƯƠNG 3: HÀM ĐẶC TRƯNG
%% ================================================================

\begin{tongket}[title={Tổng kết: Các loại hội tụ}]
\begin{enumerate}
  \item Nắm lại các \textbf{định nghĩa} và \textbf{công thức} cốt lõi trong mục ``Các loại hội tụ''.
  \item Tự làm lại ít nhất một ví dụ (nếu có) hoặc áp dụng công thức vào một tình huống số.
  \item Ghi chú điều kiện áp dụng (giả thiết độc lập, phân phối chuẩn, $n$ đủ lớn, v.v.).
\end{enumerate}
\end{tongket}

\section{Hàm đặc trưng}

\begin{dongco}[title={Động cơ --- Hàm đặc trưng}]
Công cụ gọn để chứng minh và nhận dạng phân phối qua đạo hàm tại $0$.
\end{dongco}

\begin{dinhnghia}
  $\varphi_X(t) = E[e^{itX}] = \displaystyle\int e^{itx}\,dF_X(x)$, $t \in \mathbb{R}$.
\end{dinhnghia}

\begin{tinhchat}
  \begin{itemize}
    \item $\varphi_X(0) = 1$, $|\varphi_X(t)| \le 1$, $\varphi_X$ liên tục đều.
    \item $\varphi_X$ xác định duy nhất phân phối.
    \item $X, Y$ độc lập: $\varphi_{X+Y}(t) = \varphi_X(t)\,\varphi_Y(t)$.
    \item Nếu $E[|X|^k] < \infty$: $\varphi_X^{(k)}(0) = i^k E[X^k]$.
    \item $\varphi_{aX+b}(t) = e^{ibt}\varphi_X(at)$.
  \end{itemize}
\end{tinhchat}

\begin{dinhly}[title={Định lý liên tục Lévy}]
  $X_n \xrightarrow{d} X$ $\Leftrightarrow$ $\varphi_{X_n}(t) \to \varphi_X(t)$ $\forall t$.
\end{dinhly}

\begin{vidu}
  \begin{itemize}
    \item $\mathcal{N}(\mu, \sigma^2)$: $\varphi(t) = e^{i\mu t - \sigma^2 t^2/2}$.
    \item Poisson$(\lambda)$: $\varphi(t) = e^{\lambda(e^{it}-1)}$.
    \item Exp$(\lambda)$: $\varphi(t) = \dfrac{\lambda}{\lambda - it}$.
  \end{itemize}
\end{vidu}

%% ================================================================
%%  CHƯƠNG 4: KỲ VỌNG CÓ ĐIỀU KIỆN
%% ================================================================

\begin{tongket}[title={Tổng kết: Hàm đặc trưng}]
\begin{enumerate}
  \item Nắm lại các \textbf{định nghĩa} và \textbf{công thức} cốt lõi trong mục ``Hàm đặc trưng''.
  \item Tự làm lại ít nhất một ví dụ (nếu có) hoặc áp dụng công thức vào một tình huống số.
  \item Ghi chú điều kiện áp dụng (giả thiết độc lập, phân phối chuẩn, $n$ đủ lớn, v.v.).
\end{enumerate}
\end{tongket}

\section{Kỳ vọng có điều kiện}

\begin{dongco}[title={Động cơ --- Kỳ vọng có điều kiện}]
Mô tả ``giá trị trung bình của $X$ khi đã biết thông tin $\mathcal{G}$'' --- nền của martingale.
\end{dongco}

\begin{dinhnghia}[title={Kỳ vọng có điều kiện theo $\sigma$-đại số}]
  Cho $X \in L^1$ và $\mathcal{G} \subset \mathcal{F}$ là $\sigma$-đại số con. $E[X|\mathcal{G}]$ là BNN $\mathcal{G}$-đo được duy nhất (a.s.) sao cho:
  \[
    \int_G E[X|\mathcal{G}]\,dP = \int_G X\,dP, \qquad \forall G \in \mathcal{G}.
  \]
  $E[X|Y] := E[X|\sigma(Y)]$.
\end{dinhnghia}

\begin{tinhchat}[title={Tính chất kỳ vọng có điều kiện}]
  \begin{enumerate}
    \item \textbf{Tuyến tính:} $E[\alpha X + \beta Y|\mathcal{G}] = \alpha E[X|\mathcal{G}] + \beta E[Y|\mathcal{G}]$.
    \item \textbf{Tower property:} $E\bigl[E[X|\mathcal{G}]\bigr] = E[X]$. Tổng quát: nếu $\mathcal{H} \subset \mathcal{G}$ thì $E\bigl[E[X|\mathcal{G}]\big|\mathcal{H}\bigr] = E[X|\mathcal{H}]$.
    \item \textbf{Pull-out:} Nếu $Y$ là $\mathcal{G}$-đo được thì $E[XY|\mathcal{G}] = Y\,E[X|\mathcal{G}]$.
    \item \textbf{Độc lập:} Nếu $X$ độc lập với $\mathcal{G}$ thì $E[X|\mathcal{G}] = E[X]$.
    \item \textbf{Jensen:} $\varphi$ lồi $\Rightarrow$ $\varphi(E[X|\mathcal{G}]) \le E[\varphi(X)|\mathcal{G}]$ a.s.
    \item \textbf{Hội tụ:} MCT, DCT vẫn đúng cho kỳ vọng có điều kiện.
  \end{enumerate}
\end{tinhchat}

\begin{dinhnghia}[title={Xác suất có điều kiện}]
  $P(A|\mathcal{G}) = E[\mathbf{1}_A|\mathcal{G}]$. \textbf{Phân phối có điều kiện chính quy:} tồn tại $P(B|Y = y)$ là hàm đo được theo $y$.
\end{dinhnghia}

\begin{congthuc}[title={Phương sai có điều kiện}]
  $\text{Var}(X|Y) = E[X^2|Y] - (E[X|Y])^2$.

  \textbf{Công thức phương sai toàn phần:}
  \[
    \text{Var}(X) = E[\text{Var}(X|Y)] + \text{Var}(E[X|Y]).
  \]
\end{congthuc}

%% ================================================================
%%  CHƯƠNG 5: LUẬT SỐ LỚN VÀ CLT (DẠNG ĐỘ ĐO)
%% ================================================================

\begin{tongket}[title={Tổng kết: Kỳ vọng có điều kiện}]
\begin{enumerate}
  \item Nắm lại các \textbf{định nghĩa} và \textbf{công thức} cốt lõi trong mục ``Kỳ vọng có điều kiện''.
  \item Tự làm lại ít nhất một ví dụ (nếu có) hoặc áp dụng công thức vào một tình huống số.
  \item Ghi chú điều kiện áp dụng (giả thiết độc lập, phân phối chuẩn, $n$ đủ lớn, v.v.).
\end{enumerate}
\end{tongket}

\section{Luật số lớn và CLT (góc nhìn độ đo)}

\begin{dongco}[title={Động cơ --- Luật số lớn và CLT (góc nhìn độ đo)}]
Chứng minh và điều kiện áp dụng các định lý giới hạn ở mức đo lường.
\end{dongco}

\begin{dinhly}[title={Luật mạnh số lớn (Kolmogorov)}]
  Nếu $X_n$ iid, $E[|X_1|] < \infty$ thì $\bar{X}_n \xrightarrow{\text{a.s.}} E[X_1]$.
\end{dinhly}

\begin{dinhly}[title={CLT Lindeberg--Lévy}]
  $X_n$ iid, $E[X_1] = \mu$, $\text{Var}(X_1) = \sigma^2 < \infty$:
  $\dfrac{S_n - n\mu}{\sigma\sqrt{n}} \xrightarrow{d} \mathcal{N}(0, 1)$.
\end{dinhly}

\begin{dinhly}[title={CLT Lindeberg--Feller (không đồng phân phối)}]
  Cho $X_n$ độc lập (không nhất thiết cùng phân phối), $s_n^2 = \sum_{k=1}^n \sigma_k^2$. Nếu thỏa \textbf{điều kiện Lindeberg}: $\forall \varepsilon > 0$,
  \[
    \frac{1}{s_n^2}\sum_{k=1}^n E\bigl[(X_k - \mu_k)^2\,\mathbf{1}(|X_k - \mu_k| > \varepsilon s_n)\bigr] \to 0,
  \]
  thì $\dfrac{S_n - E[S_n]}{s_n} \xrightarrow{d} \mathcal{N}(0, 1)$.
\end{dinhly}

\begin{dinhly}[title={Luật logarit lặp (Hartman--Wintner)}]
  $X_n$ iid, $E[X_1] = 0$, $\text{Var}(X_1) = 1$:
  \[
    \limsup_{n \to \infty} \frac{S_n}{\sqrt{2n\ln\ln n}} = 1 \quad\text{a.s.}
  \]
\end{dinhly}

\begin{tongket}[title={Tổng kết: Luật số lớn và CLT (góc nhìn độ đo)}]
\begin{enumerate}
  \item Nắm lại các \textbf{định nghĩa} và \textbf{công thức} cốt lõi trong mục ``Luật số lớn và CLT (góc nhìn độ đo)''.
  \item Tự làm lại ít nhất một ví dụ (nếu có) hoặc áp dụng công thức vào một tình huống số.
  \item Ghi chú điều kiện áp dụng (giả thiết độc lập, phân phối chuẩn, $n$ đủ lớn, v.v.).
\end{enumerate}
\end{tongket}



\end{document}
